\section{Results}

For running the application the following commands were used for compiling,
uploading the project and monitoring the port to observe the periodic reports.

\begin{verbatim}
# list connected boards and their ports
pio device list
# compile and upload to ESP32
pio run --target upload
# open serial monitor at 115200 baud
pio device monitor --baud 115200
\end{verbatim}

Below is the output of the commands listed above

\begin{verbatim}
======================================
  Lab 5 - Temperature Monitor
======================================
Sensor    : DS18B20 (digital) GPIO4
Thresh    : ON >= 26.0C, OFF <= 24.0C
Debounce  : 3 x 100ms = 300ms
Display   : LCD 16x2 I2C + Serial
======================================

23.5,26.0,24.0
23.6,26.0,24.0
23.7,26.0,24.0

--- SENSOR REPORT (t=500 ms) ---
Temp  : 23.5 C (valid=1, alert=OFF, dbg=0)
Thresh: ON >= 26.0 C, OFF <= 24.0 C
--------------------------------

26.1,26.0,24.0
26.2,26.0,24.0
26.3,26.0,24.0

--- SENSOR REPORT (t=5500 ms) ---
Temp  : 26.3 C (valid=1, alert=ON, dbg=0)
Thresh: ON >= 26.0 C, OFF <= 24.0 C
--------------------------------
\end{verbatim}

The system behavior is as follows:
\begin{itemize}
  \item At startup, the LCD shows \texttt{Lab5 TempMonitor} / \texttt{Starting...}
    for 1\,second while all modules initialize.
  \item During normal operation (temperature below $26\,^\circ$C), the \textbf{green LED}
    is on and the LCD line 1 shows the current temperature (e.g.\ \texttt{Temp: 23.5 C})
    while line 2 shows \texttt{Alert: OFF}.
  \item When the temperature rises above $26\,^\circ$C and the condition is confirmed for
    3 consecutive 100\,ms samples, the alert transitions to ON: the \textbf{red LED}
    turns on, the green LED turns off, and the LCD line 2 switches to \texttt{Alert: ON}.
  \item The alert turns OFF only after the temperature drops below $24\,^\circ$C and the
    lower condition is likewise confirmed for 3 consecutive samples, preventing chattering
    in the $24$--$26\,^\circ$C dead band.
  \item Every 500\,ms, a CSV line is emitted to STDIO in the format
    \texttt{temp,threshHigh,threshLow} for use with a serial plotter.
  \item Every 5\,seconds, a full structured report is printed showing the temperature
    value, sensor validity, current alert state, debounce counter, and threshold
    configuration.
\end{itemize}

Below is the visual representation of the circuit

\begin{figure}[h]
    \centering
    \includegraphics[width=0.6\textwidth]{img/ledoff.jpg}
    \caption{Circuit in idle state}
    \label{fig:ledoff}
\end{figure}

Figure \ref{fig:ledoff} shows the circuit in its idle state with both LEDs off before
acquisition begins.
